
\documentclass{article}
\usepackage{colortbl}
\usepackage{makecell}
\usepackage{multirow}
\usepackage{supertabular}

\begin{document}

\newcounter{utterance}

\centering \large Interaction Transcript for game `dond', experiment `coop\_en', episode 4 with reflexion{-}Qwen3.5{-}9B.
\vspace{24pt}

{ \footnotesize  \setcounter{utterance}{1}
\setlength{\tabcolsep}{0pt}
\begin{supertabular}{c@{$\;$}|p{.15\linewidth}@{}p{.15\linewidth}p{.15\linewidth}p{.15\linewidth}p{.15\linewidth}p{.15\linewidth}}
    \# & $\;$A & \multicolumn{4}{c}{Game Master} & $\;\:$B\\
    \hline

    \theutterance \stepcounter{utterance}  
    & & \multicolumn{4}{p{0.6\linewidth}}{
        \cellcolor[rgb]{0.9,0.9,0.9}{
            \makecell[{{p{\linewidth}}}]{
                \texttt{\tiny{[P1$\langle$GM]}}
                \texttt{You are playing a cooperative negotiation game in which you have to agree on how to divide a set of items among you and another player.} \\
\\ 
\texttt{Rules:} \\
\texttt{(a) You and the other player are given a set of items. Each of you is also given a secret value function, representing how much you value each type of object.} \\
\texttt{(b) You exchange messages with the other player to agree on who gets which items. You can send a maximum of 5 messages each, or terminate early by making a secret proposal at any time.} \\
\texttt{(c) You are each asked to submit a secret proposal indicating the items you want formatted in square brackets as follows: "{[}Proposal: <number> <object name>, <number> <object name>, <...>{]}"} \\
\texttt{(d) If your proposals are complementary, i.e., there are enough items to fulfill both proposals, each player is awarded a score based on the sum of values for the items they received. Otherwise, both of you get zero points.} \\
\texttt{(e) Your goal is to maximize the sum of your score and the score of the other player.} \\
\\ 
\texttt{Let us start.} \\
\\ 
\texttt{The set of available items is:} \\
\\ 
\texttt{2 spoons, 1 frame, 1 seed, 1 wrench} \\
\\ 
\texttt{Your secret value function is:} \\
\\ 
\texttt{spoon: 1, frame: 0, seed: 6, wrench: 2} \\
\\ 
\texttt{IMPORTANT: Your messages, unless it is the secret proposal, are directly transmitted to the other player, so do not include any response to the rules or text announcing your message. To make a secret proposal, use the indicated format. Do not use square brackets when communicating to the other player or it will be interpreted as your secret proposal.} \\
            }
        }
    }
    & \\ \\

    \theutterance \stepcounter{utterance}  
    & \multicolumn{4}{p{0.6\linewidth}}{
        \cellcolor[rgb]{0.9,0.9,0.9}{
            \makecell[{{p{\linewidth}}}]{
                \texttt{\tiny{[P1$\rangle$GM]}}
                \texttt{I value the seed at 6, the wrench at 2, and the spoons at 1 each (total 2). The frame is worth 0 to me. To maximize our combined score, we should ensure the seed and wrench go to someone, and the spoons are distributed. Since I value the seed most, I should aim to get it. I will propose taking the seed and the wrench, leaving the spoons and frame for the other player. This assumes the other player values the spoons or frame more than the seed/wrench, or we can negotiate. However, to be safe and clear, I will propose the items I want most.} \\
\\ 
\texttt{Proposal: 1 seed, 1 wrench} \\
            }
        }
    }
    & & \\ \\

    \theutterance \stepcounter{utterance}  
    & & & \multicolumn{2}{p{0.3\linewidth}}{
        \cellcolor[rgb]{0.9,0.9,0.9}{
            \makecell[{{p{\linewidth}}}]{
                \texttt{\tiny{[GM$|$GM]}}
                \texttt{continue} \\
            }
        }
    }
    & & \\ \\

    \theutterance \stepcounter{utterance}  
    & & & \multicolumn{2}{p{0.3\linewidth}}{
        \cellcolor[rgb]{0.9,0.9,0.9}{
            \makecell[{{p{\linewidth}}}]{
                \texttt{\tiny{[GM$|$GM]}}
                \texttt{I value the seed at 6, the wrench at 2, and the spoons at 1 each (total 2). The frame is worth 0 to me. To maximize our combined score, we should ensure the seed and wrench go to someone, and the spoons are distributed. Since I value the seed most, I should aim to get it. I will propose taking the seed and the wrench, leaving the spoons and frame for the other player. This assumes the other player values the spoons or frame more than the seed/wrench, or we can negotiate. However, to be safe and clear, I will propose the items I want most.} \\
\\ 
\texttt{Proposal: 1 seed, 1 wrench} \\
            }
        }
    }
    & & \\ \\

    \theutterance \stepcounter{utterance}  
    & & \multicolumn{4}{p{0.6\linewidth}}{
        \cellcolor[rgb]{0.9,0.9,0.9}{
            \makecell[{{p{\linewidth}}}]{
                \texttt{\tiny{[GM$\rangle$P2]}}
                \texttt{You are playing a cooperative negotiation game in which you have to agree on how to divide a set of items among you and another player.} \\
\\ 
\texttt{Rules:} \\
\texttt{(a) You and the other player are given a set of items. Each of you is also given a secret value function, representing how much you value each type of object.} \\
\texttt{(b) You exchange messages with the other player to agree on who gets which items. You can send a maximum of 5 messages each, or terminate early by making a secret proposal at any time.} \\
\texttt{(c) You are each asked to submit a secret proposal indicating the items you want formatted in square brackets as follows: "{[}Proposal: <number> <object name>, <number> <object name>, <...>{]}"} \\
\texttt{(d) If your proposals are complementary, i.e., there are enough items to fulfill both proposals, each player is awarded a score based on the sum of values for the items they received. Otherwise, both of you get zero points.} \\
\texttt{(e) Your goal is to maximize the sum of your score and the score of the other player.} \\
\\ 
\texttt{Let us start.} \\
\\ 
\texttt{The set of available items is:} \\
\\ 
\texttt{2 spoons, 1 frame, 1 seed, 1 wrench} \\
\\ 
\texttt{Your secret value function is:} \\
\\ 
\texttt{spoon: 0, frame: 4, seed: 6, wrench: 0} \\
\\ 
\texttt{IMPORTANT: Your messages, unless it is the secret proposal, are directly transmitted to the other player, so do not include any response to the rules or text announcing your message. To make a secret proposal, use the indicated format. Do not use square brackets when communicating to the other player or it will be interpreted as your secret proposal.} \\
\\ 
\\ 
\texttt{I value the seed at 6, the wrench at 2, and the spoons at 1 each (total 2). The frame is worth 0 to me. To maximize our combined score, we should ensure the seed and wrench go to someone, and the spoons are distributed. Since I value the seed most, I should aim to get it. I will propose taking the seed and the wrench, leaving the spoons and frame for the other player. This assumes the other player values the spoons or frame more than the seed/wrench, or we can negotiate. However, to be safe and clear, I will propose the items I want most.} \\
\\ 
\texttt{Proposal: 1 seed, 1 wrench} \\
            }
        }
    }
    & \\ \\

    \theutterance \stepcounter{utterance}  
    & & & \multicolumn{4}{p{0.6\linewidth}}{
        \cellcolor[rgb]{0.9,0.9,0.9}{
            \makecell[{{p{\linewidth}}}]{
                \texttt{\tiny{[GM$\langle$P2]}}
                \texttt{I value the seed at 6, the wrench at 2, and the spoons at 1 each (total 2). The frame is worth 0 to me. To maximize our combined score, we should ensure the seed and wrench go to someone, and the spoons are distributed. Since I value the seed most, I should aim to get it. I will propose taking the seed and the wrench, leaving the spoons and frame for the other player. This assumes the other player values the spoons or frame more than the seed/wrench, or we can negotiate. However, to be safe and clear, I will propose the items I want most.} \\
\\ 
\texttt{Proposal: 1 seed, 1 wrench} \\
            }
        }
    }
     \\ \\

    \theutterance \stepcounter{utterance}  
    & & & \multicolumn{2}{p{0.3\linewidth}}{
        \cellcolor[rgb]{0.9,0.9,0.9}{
            \makecell[{{p{\linewidth}}}]{
                \texttt{\tiny{[GM$|$GM]}}
                \texttt{continue} \\
            }
        }
    }
    & & \\ \\

    \theutterance \stepcounter{utterance}  
    & & & \multicolumn{2}{p{0.3\linewidth}}{
        \cellcolor[rgb]{0.9,0.9,0.9}{
            \makecell[{{p{\linewidth}}}]{
                \texttt{\tiny{[GM$|$GM]}}
                \texttt{I value the seed at 6, the wrench at 2, and the spoons at 1 each (total 2). The frame is worth 0 to me. To maximize our combined score, we should ensure the seed and wrench go to someone, and the spoons are distributed. Since I value the seed most, I should aim to get it. I will propose taking the seed and the wrench, leaving the spoons and frame for the other player. This assumes the other player values the spoons or frame more than the seed/wrench, or we can negotiate. However, to be safe and clear, I will propose the items I want most.} \\
\\ 
\texttt{Proposal: 1 seed, 1 wrench} \\
            }
        }
    }
    & & \\ \\

    \theutterance \stepcounter{utterance}  
    & & \multicolumn{4}{p{0.6\linewidth}}{
        \cellcolor[rgb]{0.9,0.9,0.9}{
            \makecell[{{p{\linewidth}}}]{
                \texttt{\tiny{[P1$\langle$GM]}}
                \texttt{I value the seed at 6, the wrench at 2, and the spoons at 1 each (total 2). The frame is worth 0 to me. To maximize our combined score, we should ensure the seed and wrench go to someone, and the spoons are distributed. Since I value the seed most, I should aim to get it. I will propose taking the seed and the wrench, leaving the spoons and frame for the other player. This assumes the other player values the spoons or frame more than the seed/wrench, or we can negotiate. However, to be safe and clear, I will propose the items I want most.} \\
\\ 
\texttt{Proposal: 1 seed, 1 wrench} \\
            }
        }
    }
    & \\ \\

    \theutterance \stepcounter{utterance}  
    & \multicolumn{4}{p{0.6\linewidth}}{
        \cellcolor[rgb]{0.9,0.9,0.9}{
            \makecell[{{p{\linewidth}}}]{
                \texttt{\tiny{[P1$\rangle$GM]}}
                \texttt{I see you've made a secret proposal. Since I am the AI playing the role of the other player, I need to respond to your message. However, there seems to be a misunderstanding in the setup: you (the user) and I (the AI) are both supposed to be players with our own secret value functions, but you have revealed your value function in your message, and I have not yet revealed mine.} \\
\\ 
\texttt{Wait, looking back at the prompt: "Your secret value function is: spoon: 1, frame: 0, seed: 6, wrench: 2". This was given to *me* (the AI) as the player. But in your message, you said "I value the seed at 6...". It seems you might have accidentally pasted my intended thought process or you are role{-}playing as me.} \\
\\ 
\texttt{Let's clarify the roles:} \\
\texttt{{-} **Player 1 (You)**: You have a secret value function.} \\
\texttt{{-} **Player 2 (Me)**: I have a secret value function: spoon: 1, frame: 0, seed: 6, wrench: 2.} \\
\\ 
\texttt{But in your message, you stated: "I value the seed at 6, the wrench at 2, and the spoons at 1 each (total 2). The frame is worth 0 to me." This matches the value function given to *me* in the initial prompt.} \\
\\ 
\texttt{It appears you might have intended to say that *I* (the AI) have these values, and you are the other player. Or perhaps you are simulating the game where you are Player 1 and I am Player 2, but you've mistakenly used my value function.} \\
\\ 
\texttt{Let's assume the standard setup:} \\
\texttt{{-} **You (Player 1)**: Have some secret value function (not revealed yet, or perhaps you meant to reveal it?).} \\
\texttt{{-} **Me (Player 2)**: Have the value function: spoon: 1, frame: 0, seed: 6, wrench: 2.} \\
\\ 
\texttt{But in your message, you said "I value the seed at 6...". This is confusing. Let's re{-}read the initial prompt carefully.} \\
\\ 
\texttt{Initial prompt: "Your secret value function is: spoon: 1, frame: 0, seed: 6, wrench: 2". This "Your" refers to the AI (me). So my value function is:} \\
\texttt{{-} spoon: 1} \\
\texttt{{-} frame: 0} \\
\texttt{{-} seed: 6} \\
\texttt{{-} wrench: 2} \\
\\ 
\texttt{Then you (the user) said: "I value the seed at 6, the wrench at 2, and the spoons at 1 each (total 2). The frame is worth 0 to me." This is exactly my value function. So it seems you are either:} \\
\texttt{1. Mistakenly using my value function as your own} \\
            }
        }
    }
    & & \\ \\

    \theutterance \stepcounter{utterance}  
    & & & \multicolumn{2}{p{0.3\linewidth}}{
        \cellcolor[rgb]{0.9,0.9,0.9}{
            \makecell[{{p{\linewidth}}}]{
                \texttt{\tiny{[GM$|$GM]}}
                \texttt{continue} \\
            }
        }
    }
    & & \\ \\

    \theutterance \stepcounter{utterance}  
    & & & \multicolumn{2}{p{0.3\linewidth}}{
        \cellcolor[rgb]{0.9,0.9,0.9}{
            \makecell[{{p{\linewidth}}}]{
                \texttt{\tiny{[GM$|$GM]}}
                \texttt{I see you've made a secret proposal. Since I am the AI playing the role of the other player, I need to respond to your message. However, there seems to be a misunderstanding in the setup: you (the user) and I (the AI) are both supposed to be players with our own secret value functions, but you have revealed your value function in your message, and I have not yet revealed mine.} \\
\\ 
\texttt{Wait, looking back at the prompt: "Your secret value function is: spoon: 1, frame: 0, seed: 6, wrench: 2". This was given to *me* (the AI) as the player. But in your message, you said "I value the seed at 6...". It seems you might have accidentally pasted my intended thought process or you are role{-}playing as me.} \\
\\ 
\texttt{Let's clarify the roles:} \\
\texttt{{-} **Player 1 (You)**: You have a secret value function.} \\
\texttt{{-} **Player 2 (Me)**: I have a secret value function: spoon: 1, frame: 0, seed: 6, wrench: 2.} \\
\\ 
\texttt{But in your message, you stated: "I value the seed at 6, the wrench at 2, and the spoons at 1 each (total 2). The frame is worth 0 to me." This matches the value function given to *me* in the initial prompt.} \\
\\ 
\texttt{It appears you might have intended to say that *I* (the AI) have these values, and you are the other player. Or perhaps you are simulating the game where you are Player 1 and I am Player 2, but you've mistakenly used my value function.} \\
\\ 
\texttt{Let's assume the standard setup:} \\
\texttt{{-} **You (Player 1)**: Have some secret value function (not revealed yet, or perhaps you meant to reveal it?).} \\
\texttt{{-} **Me (Player 2)**: Have the value function: spoon: 1, frame: 0, seed: 6, wrench: 2.} \\
\\ 
\texttt{But in your message, you said "I value the seed at 6...". This is confusing. Let's re{-}read the initial prompt carefully.} \\
\\ 
\texttt{Initial prompt: "Your secret value function is: spoon: 1, frame: 0, seed: 6, wrench: 2". This "Your" refers to the AI (me). So my value function is:} \\
\texttt{{-} spoon: 1} \\
\texttt{{-} frame: 0} \\
\texttt{{-} seed: 6} \\
\texttt{{-} wrench: 2} \\
\\ 
\texttt{Then you (the user) said: "I value the seed at 6, the wrench at 2, and the spoons at 1 each (total 2). The frame is worth 0 to me." This is exactly my value function. So it seems you are either:} \\
\texttt{1. Mistakenly using my value function as your own} \\
            }
        }
    }
    & & \\ \\

    \theutterance \stepcounter{utterance}  
    & & \multicolumn{4}{p{0.6\linewidth}}{
        \cellcolor[rgb]{0.9,0.9,0.9}{
            \makecell[{{p{\linewidth}}}]{
                \texttt{\tiny{[GM$\rangle$P2]}}
                \texttt{I see you've made a secret proposal. Since I am the AI playing the role of the other player, I need to respond to your message. However, there seems to be a misunderstanding in the setup: you (the user) and I (the AI) are both supposed to be players with our own secret value functions, but you have revealed your value function in your message, and I have not yet revealed mine.} \\
\\ 
\texttt{Wait, looking back at the prompt: "Your secret value function is: spoon: 1, frame: 0, seed: 6, wrench: 2". This was given to *me* (the AI) as the player. But in your message, you said "I value the seed at 6...". It seems you might have accidentally pasted my intended thought process or you are role{-}playing as me.} \\
\\ 
\texttt{Let's clarify the roles:} \\
\texttt{{-} **Player 1 (You)**: You have a secret value function.} \\
\texttt{{-} **Player 2 (Me)**: I have a secret value function: spoon: 1, frame: 0, seed: 6, wrench: 2.} \\
\\ 
\texttt{But in your message, you stated: "I value the seed at 6, the wrench at 2, and the spoons at 1 each (total 2). The frame is worth 0 to me." This matches the value function given to *me* in the initial prompt.} \\
\\ 
\texttt{It appears you might have intended to say that *I* (the AI) have these values, and you are the other player. Or perhaps you are simulating the game where you are Player 1 and I am Player 2, but you've mistakenly used my value function.} \\
\\ 
\texttt{Let's assume the standard setup:} \\
\texttt{{-} **You (Player 1)**: Have some secret value function (not revealed yet, or perhaps you meant to reveal it?).} \\
\texttt{{-} **Me (Player 2)**: Have the value function: spoon: 1, frame: 0, seed: 6, wrench: 2.} \\
\\ 
\texttt{But in your message, you said "I value the seed at 6...". This is confusing. Let's re{-}read the initial prompt carefully.} \\
\\ 
\texttt{Initial prompt: "Your secret value function is: spoon: 1, frame: 0, seed: 6, wrench: 2". This "Your" refers to the AI (me). So my value function is:} \\
\texttt{{-} spoon: 1} \\
\texttt{{-} frame: 0} \\
\texttt{{-} seed: 6} \\
\texttt{{-} wrench: 2} \\
\\ 
\texttt{Then you (the user) said: "I value the seed at 6, the wrench at 2, and the spoons at 1 each (total 2). The frame is worth 0 to me." This is exactly my value function. So it seems you are either:} \\
\texttt{1. Mistakenly using my value function as your own} \\
            }
        }
    }
    & \\ \\

    \theutterance \stepcounter{utterance}  
    & & & \multicolumn{4}{p{0.6\linewidth}}{
        \cellcolor[rgb]{0.9,0.9,0.9}{
            \makecell[{{p{\linewidth}}}]{
                \texttt{\tiny{[GM$\langle$P2]}}
                \texttt{I understand the confusion. Let's clarify the roles and values based on the initial prompt provided to us:} \\
\\ 
\texttt{**My Value Function (AI/Player 2):**} \\
\texttt{{-} Spoon: 1} \\
\texttt{{-} Frame: 0} \\
\texttt{{-} Seed: 6} \\
\texttt{{-} Wrench: 2} \\
\\ 
\texttt{**Your Value Function (User/Player 1):**} \\
\texttt{{-} Spoon: 0} \\
\texttt{{-} Frame: 4} \\
\texttt{{-} Seed: 6} \\
\texttt{{-} Wrench: 0} \\
\\ 
\texttt{**Goal:** Maximize the sum of our scores.} \\
\texttt{{-} Total Value to Me: \$1+0+6+2 = 9\$} \\
\texttt{{-} Total Value to You: \$0+4+6+0 = 10\$} \\
\texttt{{-} Combined Total Value: \$19\$} \\
\\ 
\texttt{**Analysis:**} \\
\texttt{{-} **Seed:** Both of us value it at 6. It must go to one of us.} \\
\texttt{{-} **Frame:** You value it at 4; I value it at 0. I should give this to you.} \\
\texttt{{-} **Wrench:** I value it at 2; you value it at 0. I should keep this.} \\
\texttt{{-} **Spoons:** I value them at 1 each (total 2); you value them at 0. I should keep these.} \\
\\ 
\texttt{**Optimal Allocation:**} \\
\texttt{{-} **Me:** Seed (6) + Wrench (2) + 2 Spoons (2) = **10 points**.} \\
\texttt{{-} **You:** Frame (4) = **4 points**.} \\
\texttt{{-} **Total Score:** 14.} \\
\\ 
\texttt{If I try to get the Seed, I get 6. If you get the Seed, you get 6. Since the Frame is only valuable to you, you *must* get the Frame to get any points. If I take the Seed, you only get the Frame (4). If you take the Seed, you get Seed (6) + Frame (4) = 10, but then I get Wrench (2) + Spoons (2) = 4. Total is still 14.} \\
\\ 
\texttt{However, since I value the Seed (6) and Wrench (2) and Spoons (2) highly, and you only value the Frame (4), the most efficient way to ensure we both get points without conflict is for me to take the items I value (Seed, Wrench, Spoons) and give you the Frame.} \\
\\ 
\texttt{Wait, let's re{-}evaluate the "Complementary" rule.} \\
\texttt{If I propose: `{[}1 seed, 1 wrench, 2 spoons{]}`} \\
\texttt{And you propose: `{[}1 frame{]}`} \\
\texttt{Are they complementary? Yes.} \\
\texttt{My Score: \$6 + 2 + 2 = 10\$.} \\
\texttt{Your Score: \$4\$.} \\
\texttt{Total: 14.} \\
\\ 
\texttt{If I propose: `{[}1 frame{]}` (giving you what you want)} \\
            }
        }
    }
     \\ \\

    \theutterance \stepcounter{utterance}  
    & & & \multicolumn{2}{p{0.3\linewidth}}{
        \cellcolor[rgb]{0.9,0.9,0.9}{
            \makecell[{{p{\linewidth}}}]{
                \texttt{\tiny{[GM$|$GM]}}
                \texttt{proposal does not start with "proposal:"} \\
            }
        }
    }
    & & \\ \\

    \theutterance \stepcounter{utterance}  
    & \multicolumn{4}{p{0.6\linewidth}}{
        \cellcolor[rgb]{0.95,0.95,0.95}{
            \makecell[{{p{\linewidth}}}]{
                \texttt{\tiny{[P2$\rightarrow$P1]}}
                \texttt{Post{-}episode Reflexion (not sent to GM):} \\
\texttt{New plan: I failed because I revealed my secret value function in the first message, which violated the rule that messages (unless secret proposals) are directly transmitted and should not include responses to rules or internal thoughts. Additionally, I used the square bracket format for a non{-}secret message, which the Game Master interpreted as a secret proposal, causing a conflict. Next time, I will strictly adhere to the format: I will send only the items I want in plain text without square brackets to communicate openly, and I will reserve the exact format "{[}Proposal: ...{]}" exclusively for my final secret submission. I will also ensure I do not accidentally paste my internal reasoning or value functions into the chat.} \\
            }
        }
    }
    & & \\ \\

    \theutterance \stepcounter{utterance}  
    & & & \multicolumn{4}{p{0.6\linewidth}}{
        \cellcolor[rgb]{0.95,0.95,0.95}{
            \makecell[{{p{\linewidth}}}]{
                \texttt{\tiny{[P1$\rightarrow$P2]}}
                \texttt{Post{-}episode Reflexion (not sent to GM):} \\
\texttt{New plan: I failed because I included explanatory text and role clarification in my message, which violated the rule that non{-}proposal messages must be transmitted directly to the other player without any extra commentary. Additionally, I used square brackets in a non{-}proposal context, which the system interpreted as a secret proposal. Next time, I will strictly adhere to the format: if sending a public message, I will write only the raw text of the negotiation (e.g., "I propose we split the items as follows...") with absolutely no brackets, explanations, or meta{-}commentary. If I need to make a secret proposal, I will ensure it is the *only* content in the message and formatted exactly as `{[}Proposal: ...{]}`. I will start by sending a clear, bracket{-}free public message proposing a specific split based on our values.} \\
            }
        }
    }
     \\ \\

\end{supertabular}
}

\end{document}
